%Glossary Package
\makeglossary 
%Needed to make the glossary
\usepackage{todonotes} 
%For making to-do notes in the pdf
\usepackage{float} 
%Redefine the float, now we can use 'H' instead of 'h'
\usepackage{array} 
%Extend the array environment
\usepackage{pdfpages} 
%Allows the inclusion of external pdf's 
\usepackage{color, colortbl}
%Allows us to manage the colour of items in the final output 
\usepackage{fancyhdr}
%Allows control of the headers and footers
\usepackage{nameref}
%Refer to sections by their name
\usepackage{chemfig} 
%Draw chemical diagrams
\usepackage{listings} 
%Typeset programming languages
\lstset{language=bash, numbers=left, frame=single, breaklines=true, numberstyle=\footnotesize}
%Set-up date for listings package
\usepackage{framed} 
%Framed or shaded regions that can break across pages
\usepackage[os=win]{menukeys}
%Format menu sequences, paths and keystrokes from lists
\usepackage{tikz} 
%Used for drawing within the package
\usetikzlibrary{positioning,decorations.markings,patterns}
\usepackage{multirow} 
%Create tabular cells spanning multiple rows
\usepackage[scientific-notation=true]{siunitx} 
%For putting certain scientific notation i.e. x10^.
\usepackage{titlesec} 
%Custom titles
\usepackage[normalem]{ulem}
%Strikethrough text with \sout{}
\usepackage{titlesec}
\usepackage{tikz}
\renewcommand{\familydefault}{\sfdefault}

\titleformat*{\section}{\Huge\bfseries}
\titleformat*{\subsection}{\LARGE\bfseries}
\titleformat*{\subsubsection}{\Large\bfseries}

\setcounter{tocdepth}{3}
\setcounter{secnumdepth}{3}